\documentclass[11pt,a4paper]{article}
\usepackage[utf8]{inputenc}
\usepackage[T1]{fontenc}
\usepackage[brazil]{babel}
\usepackage{lmodern}
\usepackage[margin=2.1cm]{geometry}
\usepackage{microtype,booktabs,tabularx,enumitem}
\usepackage[hidelinks]{hyperref}
\setlength{\parindent}{0pt}
\setlength{\parskip}{5pt}
\setlist{nosep,leftmargin=*}
\newcommand{\msg}[1]{\texttt{\detokenize{#1}}}
\title{Chat com ordenação\\\large Documento de protocolo -- esboço v0}
\author{Levi \and Henri \and André \and Gabriel}
\date{29 de outubro de 2026}
\begin{document}
\maketitle
\vspace{-1.4em}

\textbf{Estado.} Esboço revisável, sem implementação nem validação em execução.
Os campos e comportamentos abaixo orientam a primeira versão; decisões abertas
devem ser resolvidas antes de tratar o protocolo como contrato definitivo.
Arquitetura e justificativas ficam em \msg{docs/relatorio.tex}.

\section{Escopo e transporte}
Três nós equivalentes, identificados por 1, 2 e 3, trocam JSON por HTTP.
Clientes usam WebSocket. Nomes dos serviços vêm da configuração e são resolvidos
pelo DNS do Compose; não há IPs fixos no código.
O cenário inicial é uma sala, com estado em memória.

A recuperação proposta pressupõe uma queda por vez, rede funcionando entre os
sobreviventes e nenhuma nova falha durante reorganização ou sincronização.
São necessários dois dos três nós para operar. O relatório detalha os limites.

\section{Protocolo v0: operação normal}
\textbf{Identidade e visão.} Cada nó tem ID 1, 2 ou 3.
Uma visão contém época, líder e membros ativos. A época é proposta como o par
(contador, ID do coordenador), que distingue propostas sem substituir acordo
entre participantes. Estados locais: sincronizando, ativo e reorganizando.

Cada nó mantém relógio $L$, fila de pendências, ACKs recebidos, histórico,
UUIDs entregues e identificação dos pacotes já processados.
O histórico preserva o conteúdo, a época e o carimbo original de cada evento,
necessários para a reconciliação; guardar apenas UUIDs não basta.

\textbf{Cliente.} Envia \msg{id}, \msg{sala}, \msg{autor} e \msg{texto};
\msg{id} é um UUID preservado nos reenvios.
Reutilizar um UUID com conteúdo diferente é erro de protocolo.

\textbf{Envelope entre nós.} Campos comuns: \msg{tipo}, \msg{pacote_id},
\msg{epoca} e \msg{origem}. Um retry preserva o mesmo pacote.
\msg{MSG} acrescenta relógio e conteúdo da mensagem; \msg{ACK} acrescenta
seu relógio e a referência ao evento reconhecido.
O evento é identificado pela época, origem e timestamp original.

\begin{enumerate}
\item Ao aceitar uma mensagem do cliente, o nó incrementa $L$, registra o evento
e enfileira seu multicast para todos os membros, incluindo recepção local.
\item Ao receber \msg{MSG}, atualiza $L \leftarrow \max(L,t)+1$ e guarda a
mensagem na fila, ordenada pelo \textbf{carimbo original} $(t,\mathrm{origem})$.
O relógio atualizado do receptor não altera essa chave.
\item O receptor incrementa seu relógio para emitir \msg{ACK} a todos os membros.
Receber um ACK também atualiza $L$ pela regra de recepção.
ACKs que chegam antes do MSG correspondente são guardados.
\item A primeira mensagem da fila só é entregue após ACK de todos os membros
daquela visão. O UUID é registrado no histórico e os clientes locais são avisados.
\end{enumerate}

\textbf{Confirmação.} O nó do cliente confirma a mensagem quando a entrega.
Receber pelo WebSocket ou responder a um POST não confirma entrega.
Na reconexão, o cliente consulta ou reenvia usando o mesmo UUID.
Os detalhes desse contrato com o navegador ainda serão fechados.

\textbf{Relógio.} Modela os eventos da camada \msg{MSG}/\msg{ACK}.
Heartbeats ficam fora dessa ordenação. Dentro de uma visão, a chave é
$(t,\mathrm{origem})$; entre visões, o histórico anterior deve permanecer
como prefixo imutável. Lamport ordena mensagens; o líder coordena a visão.

\subsection*{FIFO, retries e deduplicação}
Uma fila de saída por destino envia um POST por vez e só avança após a resposta.
Falhas de transporte repetem o mesmo pacote. Uma resposta pode se perder depois
do processamento: o receptor deve responder ao retry sem repetir seus efeitos.

Atualização de estado e enfileiramento de cada multicast são serializados.
O multicast inteiro é enfileirado antes de processar outro evento.
O handler não espera respostas de outros nós para concluir o processamento local.

A deduplicação ocorre no transporte e na entrega. Pacotes repetidos não geram
novo evento nem nova emissão de ACK; um UUID entregue não reaparece no histórico.
A identidade dos pacotes após reinício e a tentativa simultânea por nós distintos
precisam de detalhamento antes da implementação.

\subsection*{Heartbeat e eleição}
Parâmetros iniciais: heartbeat a cada 1 segundo e suspeita após 3 segundos.
São configuráveis e não representam prazo garantido de recuperação.
Com líder disponível, ele coordena a reorganização. Se ele for suspeito, os
membros ativos executam Bully com \msg{ELECTION}, \msg{OK} e \msg{COORDINATOR}.
Nós sincronizando ficam fora da eleição. Vencer não libera envios automaticamente.

\section{Reorganização e reentrada: proposta a revisar}
A transição fecha a visão antiga antes de retomar a aplicação.
A sequência abaixo é um esboço de recuperação sob as hipóteses da seção 1;
sua segurança ainda deve ser examinada em exemplos e testes.

\textbf{O que significa o corte.} Visão é a configuração compartilhada de membros,
líder e época. O corte é a pausa local que separa o processamento normal da visão
antiga e a coleta do estado para reorganização. Não é um horário global nem um
limite no relógio de Lamport: cada nó pausa em um momento diferente.
A pausa e a captura de estado são serializadas com o processamento local.
Uma mensagem aceita antes da pausa aparece no histórico ou nas pendências
informadas; depois da pausa, MSGs e ACKs antigos deixam de alterar esse estado.

\begin{enumerate}
\item \msg{VIEW_PREPARE}: o coordenador informa a transição e os membros propostos.
O próprio coordenador também executa a pausa e participa da coleta.
\item Cada participante pausa novos envios e entregas, encerra retries de
\msg{MSG}/\msg{ACK} da visão antiga e captura seu estado. Cópias antigas desses
dois tipos recebidas depois da pausa são descartadas com resposta explícita,
encerrando o retry. Os controles válidos da transição continuam sendo processados;
a regra de descarte não pode bloquear a reorganização.
\item \msg{VIEW_STATE}: todos os membros propostos informam histórico,
pendências e relógio. O coordenador aguarda todos.
\item O coordenador reúne os eventos conhecidos, ordenando os da visão encerrada
pelos carimbos originais. O resultado deve preservar, como prefixo, cada histórico
já entregue. Incompatibilidade mantém a pausa e sinaliza erro.
\item \msg{VIEW_INSTALL}: distribui nova época, membros, histórico reconciliado
e relógio de referência, ao menos o maior relógio informado na coleta.
Os nós não reduzem seu relógio local. Instalam o mesmo estado e permanecem pausados,
preservando o sufixo ainda não notificado aos seus clientes.
\item \msg{VIEW_READY}: cada membro confirma a instalação.
Só após todos confirmarem, o coordenador envia \msg{VIEW_RESUME}.
Na retomada, cada nó notifica seus clientes somente das mensagens que faltavam
em seu histórico anterior, preservando a ordem e os UUIDs já entregues.
\end{enumerate}

Mensagens da nova visão recebidas antes da retomada aguardam processamento.
Pacotes antigos não retornam à operação normal.
Pendências conhecidas por sobreviventes podem ser recuperadas. Uma tentativa
não confirmada conhecida apenas pelo nó que caiu pode desaparecer e exigir reenvio.
Uma nova falha durante essa transição fica fora da garantia inicial.

\subsection*{Exemplo revisado: mensagens em trânsito}
A visão inicial é época $(1,3)$, membros $\{1,2,3\}$ e líder 3.
Os históricos já contêm a mensagem \msg{oi}.
O nó 1 origina \msg{x}, com carimbo $(10,1)$; o nó 3 origina \msg{y},
com carimbo $(20,3)$, e cai depois de enviar cópias.

Antes das pausas, o nó 1 tem todos os ACKs de \msg{x} e a entrega.
O nó 2 ainda espera o ACK do nó 1 para \msg{x}; já recebeu o ACK do nó 3.
A cópia de \msg{y} destinada ao nó 1 chegou; a destinada ao nó 2 está em trânsito.
Esse cenário respeita FIFO por par de nós: o ACK do nó 3 para \msg{x}
chegou ao nó 2 antes de sua cópia de \msg{y}.

\begin{tabularx}{\linewidth}{@{}lXX@{}}
\toprule
Estado após a pausa local & Nó 1 & Nó 2 \\
\midrule
Histórico entregue & \msg{[oi, x]} & \msg{[oi]} \\
Pendências conhecidas & \msg{[y]} & \msg{[x]} \\
\bottomrule
\end{tabularx}

O Bully escolhe o nó 2 entre os sobreviventes. Ambos pausam e informam os estados
acima. Se a cópia antiga de \msg{y} chegar ao nó 2 depois de sua pausa, será
descartada nesse caminho normal. Seu conteúdo, porém, já foi preservado no relato
do nó 1 e será recuperado pela transição.

A união ordenada dos relatos produz \msg{[oi, x, y]}, que estende os dois
históricos. A proposta instala a visão de época $(2,2)$, membros $\{1,2\}$,
líder 2 e esse histórico. Após os dois READYs, a retomada entrega apenas
\msg{y} aos clientes do nó 1 e \msg{x, y} aos do nó 2.
As pendências recuperadas são resolvidas pela instalação comum, sem esperar ACK
do membro removido. A instalação pode recuperar uma mensagem que ainda não tinha
sido confirmada; ausência de confirmação não significa cancelamento.

Uma mensagem \msg{z} conhecida somente pelo nó 3 que caiu pode não aparecer
na coleta. Ela não poderia ter sido confirmada pelo fluxo normal sem ACK dos
demais membros. Seu cliente poderá repetir a tentativa com o mesmo UUID.
Este exemplo não cobre falhas adicionais durante a reorganização.

\subsection*{Reentrada}
O nó retorna sincronizando e envia \msg{JOIN}.
Sua admissão usa a mesma pausa. Recebe \msg{SNAPSHOT} com histórico, UUIDs
entregues, visão e relógio antes de entrar no grupo ativo.
Depois pode participar de uma eleição Bully, inclusive se tiver o maior ID.


\section{Mensagens de controle}
Os tipos abaixo também usam o envelope comum. O payload exato, a validação
dos campos e as respostas HTTP ainda serão definidos.

\begin{tabularx}{\linewidth}{@{}lX@{}}
\toprule
Tipo & Conteúdo ou função pretendida \\
\midrule
\msg{HEARTBEAT} & Sinal de vida; campos de estado a fechar. \\
\msg{ELECTION} / \msg{OK} & Iniciar eleição e responder como candidato de ID maior. \\
\msg{COORDINATOR} & Anunciar o vencedor, sem liberar envios da aplicação. \\
\msg{VIEW_PREPARE} & Identificar a transição e os membros propostos. \\
\msg{VIEW_STATE} & Informar histórico, pendências e relógio após o corte. \\
\msg{VIEW_INSTALL} & Distribuir época, membros, histórico e relógio de referência. \\
\msg{VIEW_READY} & Confirmar a instalação da transição correspondente. \\
\msg{VIEW_RESUME} & Liberar a operação após a confirmação de todos. \\
\msg{JOIN} & Solicitar admissão de um nó em sincronização. \\
\msg{SNAPSHOT} & Transferir histórico, UUIDs entregues, visão e relógio. \\
\bottomrule
\end{tabularx}

\section{Decisões de contrato ainda abertas}
Fechar o bootstrap da primeira visão, os tipos e nomes exatos dos campos,
as respostas HTTP e a referência de transição usada nas mensagens de controle.
Detalhar a identidade após reinício, o tratamento do mesmo UUID recebido por nós
distintos, a implementação da pausa local e a reconciliação de pendências.
Definir também como o navegador consulta, confirma e reenvia tentativas.

A revisão deve simular mensagens em trânsito e conferir a preservação do prefixo
já entregue. O exemplo acima foi conferido com um modelo de dados em memória:
união ordenada, prefixos, sufixos a entregar e rejeição de prefixo incompatível.
Isso não executou nós, transporte ou o protocolo e não prova sua segurança geral.
\section*{Referências iniciais}
Enunciado do professor: Trabalho Final -- Sistemas Distribuídos 2026,
fornecido durante o planejamento.
\href{https://lamport.azurewebsites.net/pubs/time-clocks.pdf}{Lamport (1978)};
\href{https://homepage.cs.uiowa.edu/~ghosh/Bully.pdf}{Garcia-Molina (1982)};
\href{https://www.cs.cornell.edu/home/rvr/papers/podc/node7.html}{Van Renesse et al.,
comunicação e mudança de visão}.
\end{document}
